\part*{РЕФЕРАТ}
\addcontentsline{toc}{part}{РЕФЕРАТ}

Расчетно-пояснительная записка к выпускной квалификационной работе «Метод обнаружения фальсифицированного фрагмента на изображении с использованием сверточной нейронной сети» содержит 61 страниц, 4 части, \totalfigures\ рисунков,  \totaltables\ таблицы, список используемых источников из 34 наименований и 2 приложения.

Ключевые слова: компьютерное зрение, глубокое обучение, сверточная нейронная сеть.

Объект разработки – метод обнаружения фальсифицированного фрагмента на изображении.

Цель работы: разработка и программная реализация метода автоматизированного определения области фальсифицированного фрагмента изображения с помощью двухпоточной свёрточной нейронной сети.

В первой части работы приведены обзор и сравнение существующих методов обнаружения фальсифицированного фрагмента на изображении.

Во второй части описан метод обнаружения фальсифицированного фрагмента на изображении с использованием сверточной нейронной сети, приведены IDEF0-диаграммы метода и схемы его алгоритмов.

В третьей части обоснован выбор средств разработки, приведена структура ПО, реализующий предложенный метод и описано взаимодействие пользователя с программным обеспечением..

В четвертой части выбраны метрики для оценки качества работы метода и приведены результаты исследования их зависимости от конфигурации нейронной сети и объемов передаваемых данных.

Разработанный метод может найти применение в области информационной безопасности и верификации медиаконтента в средствах массовой информации и социальных сетях. Это позволит своевременно выявлять и блокировать распространение дезинформации.
